\hypertarget{the-trusted-base}{%
\chapter{4. The trusted base}\label{the-trusted-base}}

Chapter 3 ended on a single crate's division of labour: the unsafe
binder ABI beneath the policy brain that holds none of it. That division
scales to the whole codebase, and what it scales to is the trusted base,
the set of crates whose correctness the confinement rests on. How large
that base is, how it is held small, and where the honest count of it
parts from the flattering one are what follow.

\hypertarget{what-counts-as-trusted}{%
\section{4.1 What counts as trusted}\label{what-counts-as-trusted}}

The trusted computing base is the code linked into the binaries that
enforce: \texttt{kenneld} the daemon, \texttt{kennel-privhelper} the
capability helper, and \texttt{kennel-bin-init} the root-owned PID 1. A
bug in any of them can break confinement, so their first-party
dependency closure is the base whose correctness is load-bearing. At the
time of writing, between releases 0.4.0 and 0.5.0, the whole workspace
is 27 crates and about 35,000 lines; the TCB closure is 16 of them,
about 21,000. The other 11, some 13,600 lines, are outside it: the
operator CLI and the crates only it pulls, the in-kennel facades the
workload reaches, and the D-Bus mediation engine.

What keeps those 11 out is a crate boundary, not a habit of not calling
them. The operator CLI, the CLI-to-daemon control protocol, and the
policy compiler are each their own crate, so the daemon links only what
it needs to verify a settled policy and act on it. The compiler is the
clearest case: it is the larger half of the policy code, about 5,000
lines of source schema, template resolution, deltas, signing, lint, and
risk evaluation, and the daemon links none of it. Authoring depends on
the runtime half and never the reverse, and only the CLI links the
compiler, so a \texttt{cargo\ tree} over the daemon shows no compiler,
no argument parser, no JSON. The minimisation is structural and
checkable, not a property the daemon happens to have today
(\texttt{rule-of-1}).

\hypertarget{the-vendored-truth}{%
\section{4.2 The vendored truth}\label{the-vendored-truth}}

The lines counted so far are first-party, and that is the flattering
count. The base a compromise runs through is dominated by the vendored
crates the trusted crates pull in: at the same snapshot, roughly 200,000
lines of vendored code against the 21,000 first-party in the closure, an
order of magnitude more. An inventory that stopped at first-party would
understate the audit surface by that order of magnitude, so the honest
one counts the vendored code too, and sorts it on two axes: whether it
is logic or bindings, and whether the input it sees is adversarial or
trusted (\texttt{own-your-supply-chain}, \texttt{read-by-the-hostile}).

Most of the vendored bulk is bindings, not logic. \texttt{libc} alone is
over 100,000 lines, and it is almost entirely \texttt{extern}
declarations and constants; \texttt{nix} is another 30-odd thousand of
typed wrappers over the same surface. These are platform declarations
resolving to the system's own \texttt{libc.so} and the kernel,
cfg-gated, with close to no per-line algorithmic risk, because nothing
in them runs an algorithm over an input. The logic that does run
in-process is far smaller, on the order of 60,000 vendored lines, and it
is a short list: an ELF reader, the serialisation of the project's own
typed structures, a signature verifier, the compiler for the project's
own seccomp filter, and the parser for the signed policy artefact.

\hypertarget{the-empty-intersection}{%
\section{4.3 The empty intersection}\label{the-empty-intersection}}

Sorting on two axes is worth the trouble because of the one cell it
leaves empty. Logic against bindings says what executes in the process;
adversarial against trusted says what an attacker can steer. The cell
that would hurt, vendored logic running on adversarial input inside the
daemon, has nothing in it. The vendored logic that runs in the daemon
all reads trusted input: the project's own ELF, its own typed wire, its
own filter program, the signed and verified policy artefact. The one
piece that touches attacker-supplyable bytes is the signature verifier,
and verifying those bytes is the whole point of it, a guard rather than
a victim. The ANSI terminal parser that does read workload-controlled
output, the obvious place adversarial vendored logic would creep in,
runs in the operator CLI and not at the daemon's read point, where the
broker is a raw-byte router that parses nothing.

The one first-party parser of adversarial bytes left in the daemon is
the binder command decoder, exactly the one chapter 3 quarantined and
fuzzed (T5.1). Everything else adversarial is pushed to the untrusted
side by construction. The D-Bus engine makes the discipline visible: it
parses the hostile D-Bus wire, some 4,000 lines of it, but it is
vendored onto the in-kennel facade and the operator-context delegate,
never onto the daemon, so a \texttt{cargo\ tree} over the daemon stays
free of it. Mediating D-Bus grew the facade's attack surface, on the
untrusted side of the boundary, and left the trusted base unchanged. The
host-side delegates are the same move for I/O rather than parsing: the
outbound dialer and the inbound accepter are dumb processes, each close
to a netcat over an owner-only socket, while the daemon holds all the
policy and none of the blocking work (\texttt{control-not-data-plane}).

\hypertarget{the-unsafe-quarantined}{%
\section{4.4 The unsafe, quarantined}\label{the-unsafe-quarantined}}

Five crates carry \texttt{unsafe}, and they are the same five throughout
the system: the raw-syscall, namespace, and seccomp surface; the
hand-rolled Landlock ABI; the BPF loader; the binder ioctl ABI; and the
\texttt{SCM\_RIGHTS} descriptor adoption. Together they are under 4,000
lines, and all of them are leaves or near-leaves in the dependency
graph, depended on rather than depending. The safe operating-system
primitives, path canonicalisation, uid and gid handling, the netlink and
userns-map handshakes, were deliberately split out into their own safe
crate so the unsafe crates stay small enough to read in one sitting,
which is the size a reviewer can hold in their head and a fuzzer can
cover (\texttt{quarantine-the-unsafe}, \texttt{rule-of-1}). The trusted
base is large in total, most of it vendored bindings that run no
algorithm, and the part of it that is both first-party and unsafe is
under 4,000 lines a reviewer can read end to end.

\hypertarget{the-shape-of-the-workspace}{%
\section{4.5 The shape of the
workspace}\label{the-shape-of-the-workspace}}

The naming makes the inventory legible. A \texttt{kennel-lib-} crate is
a library; a \texttt{host-} binary is a host-side delegate running in
the operator's context, outside every kennel; a \texttt{facade-} binary
is the in-kennel end of a connector, confined workload-side code; the
rest are the daemon, the helper, and the inits. The dependency graph has
a direction: \texttt{kenneld} sits at the top, pulling nearly every
library to do its work, the unsafe mechanism crates sit near the bottom
as leaves providing one primitive each, the libraries depend downward,
and the authoring half depends on the runtime half and never the
reverse, so there are no cycles and the trusted closure is exactly the
downward cone under the three enforcing binaries. Of those three only
one holds host privilege, \texttt{kennel-privhelper}; \texttt{kenneld}
runs unprivileged as the operator and \texttt{kennel-bin-init} holds no
host capability, as chapter 2 set out. The trusted base is defined by
what would break confinement, not by what holds privilege, and most of
it holds none. The whole workspace is blocking and thread-per-connection
with no async runtime, which keeps the call graph something a reader can
follow by hand.
